\chapter{电力电子器件应用的共性问题}

第 2 章讲了器件本身，本章讲器件\textbf{用到电路里}时的共性问题：怎么驱动它开通与关断、怎么保护它不被过电压过电流损坏、电压或电流不够时怎么串并联使用。这三个问题是所有电力电子装置都要面对的工程问题，与具体电路拓扑无关。

% ============================================================
\section{电力电子器件的驱动}

\subsection{电力电子器件驱动电路概述}

\textbf{驱动电路}是主电路与控制电路之间的接口，其作用是：把信息电子电路发出的控制信号\textbf{放大}到足以使器件可靠开通或关断的功率水平，并实现主电路与控制电路之间的\textbf{电气隔离}。

\begin{itemize}
  \item \textbf{隔离方式}：主电路为高压、控制电路为低压，必须隔离。常用\textbf{光耦合器}（普通型、高速型、高传输比型）与\textbf{脉冲变压器}（磁隔离）。光耦的响应时间可做到 $1.5\,\mu$s 以下；脉冲变压器在脉冲较宽时须采用\textbf{高频调制与解调}，以避免铁心饱和。
  \item \textbf{驱动电路的性能直接决定器件的开关特性}：驱动电流/电压的幅值、前沿陡度、上升下降时间，影响开通关断损耗与开关时间；驱动电路还应具备\textbf{保护功能}（如过流时软关断、欠压时封锁）。
  \item 工程上应优先选用器件生产厂家配套的\textbf{专用驱动电路/驱动模块}，以获得参数的最佳配合。
\end{itemize}

驱动电路按控制信号性质分为\textbf{电流驱动型}（晶闸管、GTO、GTR）与\textbf{电压驱动型}（电力 MOSFET、IGBT）两大类。

\subsection{晶闸管的触发电路}

晶闸管的驱动电路常称为\textbf{触发电路}，其作用是产生符合要求的门极触发脉冲，保证晶闸管在需要的时刻由阻断转为导通。对触发电路的要求（也是第 3 章 3.8 节的要求）：

\begin{enumerate}
  \item 触发脉冲应具有足够的\textbf{幅值和功率}（一般门极电压 $4\sim10$ V、电流几百毫安到几安）；
  \item 脉冲\textbf{前沿要陡}（上升时间短），以提高触发时刻的准确性、避免局部过热；
  \item 脉冲\textbf{宽度足够}（阻感负载需宽脉冲或脉冲列），保证阳极电流上升到擎住电流以上；
  \item 具有足够的\textbf{移相范围}，并与主电路电源\textbf{同步}；
  \item 各相脉冲应\textbf{对称}（间隔均匀），否则会造成输出电压不对称、产生直流分量。
\end{enumerate}

常用实现方式：阻容移相桥（简单但精度低）、单结晶体管触发电路（简单、用于小功率）、\textbf{锯齿波同步触发电路}（应用最广）、集成触发器（KJ004 系列）、数字触发器（精度可达 $0.7^\circ$ 以内）。

\subsection{典型全控型器件的驱动电路}

\subsubsection*{1. 电流驱动型器件的驱动电路}

\textbf{GTO 的驱动}：GTO 需要\textbf{正、负两个电源}——开通时加正向门极电流（幅值较小、持续整个导通期），关断时加\textbf{幅值很大的负门极电流脉冲}（因电流关断增益 $\beta_{off}$ 只有 5 左右，关断电流 $I_{GM}\approx I_{ATO}/5$）。因此 GTO 驱动电路复杂、功耗大，须用强脉冲加前沿陡化措施。

\textbf{GTR 的驱动}：基极驱动电流应使 GTR 工作在\textbf{准饱和状态}（防止过饱和导致关断时间变长），常用\textbf{抗饱和（贝克箝位）电路}；关断时须加\textbf{反向基极电流}以加速存储电荷抽出。GTR 驱动电路还需具备过流保护（退饱和检测）。

\subsubsection*{2. 电压驱动型器件的驱动电路}

\textbf{电力 MOSFET 的驱动}：本质是给栅极电容充放电。要求：
\begin{itemize}
  \item 开通时加足够幅值的正栅压（一般 $10\sim15$ V），保证充分导通、$R_{on}$ 最小；
  \item 关断时加\textbf{负栅压}（$-5\sim-15$ V），防止 du/dt 引起的误导通；
  \item 串接\textbf{栅极电阻} $R_g$ 控制充放电速度：$R_g$ 小则开关快、损耗小但 EMI 大；$R_g$ 大则相反；
  \item 驱动电路应能提供足够的峰值电流（$\pm$ 数安），并具有隔离。
\end{itemize}

\textbf{IGBT 的驱动}：与 MOSFET 类似但要求更高——栅极电荷大、有\textbf{米勒效应}（关断时集电极电压上升通过米勒电容耦合到栅极，可能引起误导通）。典型要求：正栅压 $+15$ V、负栅压 $-5\sim-15$ V；具备\textbf{退饱和（$U_{CE(sat)}$）检测}以实现短路保护，并采用\textbf{软关断}（缓慢降低栅压）限制关断过电压。常用集成驱动模块（如 EXB841、M57962 等）把这些功能集成在一起。

\corepoint{驱动电路的三句话：\textbf{（1） 晶闸管要“脉冲”（够宽够陡），GTO/GTR 要“正负双电源、强脉冲”（电流驱动），MOSFET/IGBT 要“正负栅压 + 栅极电阻”（电压驱动）}；\textbf{（2） 隔离方式看场合：光耦最常用，脉冲变压器需高频调制}；\textbf{（3） 驱动电路不只是“放大”，它还承担保护（退饱和检测、软关断、欠压封锁）}。}

% ============================================================
\section{电力电子器件的保护}

\subsection{过电压的产生及过电压保护}

电力电子装置中可能发生的过电压分为\textbf{外因过电压}与\textbf{内因过电压}两类。

\begin{table}[htbp]
\centering
\caption{过电压的分类与保护措施}
\label{tab:9-1}
\begin{tabular}{p{2.6cm}p{4.6cm}p{6.4cm}}
\toprule
\textbf{类别} & \textbf{产生原因} & \textbf{保护措施} \\
\midrule
外因过电压（来自外部） & 雷击过电压；供电系统操作过电压（断路器分合、熔断器熔断引起的电网电压波动） & 装置入口加\textbf{避雷器、压敏电阻、硒堆}；变压器屏蔽与静电屏蔽；进线加 RC 吸收 \\
内因过电压（装置内部） & \textbf{换相过电压}（器件关断时回路电感储能释放）；\textbf{关断过电压}（di/dt 大、线路电感产生尖峰）；器件反向恢复引起的尖峰 & 器件两端加 \textbf{RC 吸收电路}、\textbf{RCD 钳位}；合理安排布线以减小线路电感；选用软恢复二极管 \\
\bottomrule
\end{tabular}
\end{table}

\subsection{过电流保护}

电力电子电路运行不正常或发生故障时可能发生过电流，其原因有：\textbf{负载过载}、\textbf{直流侧短路}、\textbf{逆变失败（逆变颠覆）}、\textbf{器件损坏造成直通}等。过电流保护手段及其配合：

\begin{enumerate}
  \item \textbf{快速熔断器}：直接串在器件支路中，动作最快，用于短路保护；但熔断后须更换，且熔断过程中器件仍承受短路电流冲击。
  \item \textbf{过电流继电器 / 直流快速断路器}：用于负载过载与较大故障电流，动作较慢。
  \item \textbf{电子保护（电流检测 + 封锁脉冲）}：用霍尔电流传感器或分流器检测电流，超过阈值时立即\textbf{封锁触发脉冲}（或关断全控器件），动作最快、不损坏器件，是装置的主要保护手段；缺点是须与器件电流容量配合整定。
  \item \textbf{过载保护（热保护）}：按器件结温或散热器温度动作，用于长期过载。
\end{enumerate}

保护配合的原则是\textbf{“分级配合、由快到慢”}：电子保护最先动作（微秒级），快速熔断器作为后备（毫秒级），断路器最后（几十毫秒到秒级）。

\subsection{缓冲电路}

\textbf{缓冲电路}（Snubber，又称吸收电路）的作用是：抑制器件的 du/dt、di/dt 与电压电流尖峰，\textbf{改变器件的开关轨迹，使之落在安全工作区（SOA）内}，并可减少开关损耗（把损耗从器件转移到缓冲元件上）。

\begin{itemize}
  \item \textbf{关断缓冲电路（RC / RCD）}：并联在器件两端，器件关断时电容吸收回路电感释放的能量、限制 du/dt；开通时电容经电阻放电（RCD 中用二极管旁路电阻，使放电加快、减少开通损耗）。
  \item \textbf{开通缓冲电路}：串联电感限制开通时的 di/dt，避免电流集中与反向恢复尖峰。
  \item \textbf{复合缓冲电路}：同时具有开通与关断缓冲功能（如 $L$、$C$、$D$ 组合）。
  \item \textbf{无损缓冲电路}：用电感、电容、二极管构成，缓冲能量回馈到负载或电源而不消耗在电阻上，效率高，但结构复杂。
\end{itemize}

\corepoint{保护这一节的考点是“对号入座”：\textbf{外因过电压→避雷器/压敏电阻；内因过电压→RC 吸收、RCD 钳位；短路→快速熔断器；过载→电子保护/热保护；du/dt 过大→RC 缓冲；di/dt 过大→串联电感；二次击穿→GTR 的 SOA 与缓冲电路}。缓冲电路的核心作用是\textbf{把开关轨迹限制在 SOA 内}。}

% ============================================================
\section{电力电子器件的串联和并联使用}

当单个器件的电压或电流定额不能满足要求时，需将器件\textbf{串联}以提高耐压、\textbf{并联}以提高通流能力。但器件特性不可能完全一致，必须解决\textbf{均压}与\textbf{均流}问题。

\subsection{晶闸管的串联}

晶闸管串联时，由于各器件\textbf{反向漏电流与正反向阻断特性不同}，承受的电压不相等，须采取均压措施：

\begin{itemize}
  \item \textbf{静态均压}：在每只器件两端并联\textbf{阻值相等的均压电阻}，使静态电压按电阻分压而均匀；
  \item \textbf{动态均压}：在每只器件两端并联 \textbf{RC 吸收电路}，利用电容使开通关断过程中的电压变化均匀（动态均压电容同时起吸收作用）；
  \item \textbf{触发同时性}：串联器件必须\textbf{同时触发}，否则先导通的器件承受全部电压而损坏——须用强触发、同一触发信号源。
\end{itemize}

串联器件数按下式估算（考虑均压系数 $K_u$）：

\keyformula{n\ge\frac{U_m}{K_u\,U_{RRM}},\qquad K_u=0.8\sim0.9}

\subsection{晶闸管的并联}

晶闸管并联时，由于各器件\textbf{正向伏安特性不同}，同一电压下流过的电流不相等，须采取均流措施：

\begin{itemize}
  \item \textbf{串联电阻均流}：在每只器件支路串入小电阻，用电阻的压降强迫电流均匀（简单但损耗大）；
  \item \textbf{串联电抗器均流}：串入铁心电抗器（或空心电抗器），利用电感对电流变化的阻碍实现动态均流；
  \item \textbf{选配器件}：挑选\textbf{同一型号、同一批次、正向压降相近}的器件并联；
  \item \textbf{降额使用}：并联后总电流定额按均流系数 $K_i=0.8\sim0.9$ 折算。
\end{itemize}

\keyformula{n\ge\frac{I_m}{K_i\,I_{T(AV)}},\qquad K_i=0.8\sim0.9}

\subsection{电力 MOSFET 的并联和 IGBT 的并联}

电力 MOSFET 的通态电阻 $R_{on}$ 具有\textbf{正温度系数}：某只管子电流偏大时温度升高、$R_{on}$ 增大，电流自动减小——这种\textbf{自动均流}特性使 MOSFET 的并联比晶闸管容易得多。

IGBT 的通态压降 $U_{CE(sat)}$ 也有一定的正温度系数（在较大电流下），故也可并联，但仍需注意：

\begin{itemize}
  \item 各管\textbf{阈值电压 $U_{GE(th)}$ 与跨导 $g_{fs}$ 的差异}会造成开通时刻与电流分配不一致，须\textbf{各自串接栅极电阻}（不要共用一只电阻），必要时单独整定；
  \item 并联时须\textbf{对称布线}，使各支路的引线电感尽量相等，否则动态电流分配不均；
  \item 一般按 \textbf{90\%} 降额使用，并保证散热条件一致（同一散热器、对称安装）。
\end{itemize}

\corepoint{本节抓三条：\textbf{（1） 晶闸管串联靠“均压电阻 + RC 动态均压 + 同时触发”，并联靠“串电阻/电抗 + 选配 + 降额”}；\textbf{（2） 均压系数、均流系数一般取 0.8$\sim$0.9，串联/并联数按此折算}；\textbf{（3） MOSFET 因 $R_{on}$ 正温度系数可自动均流，IGBT 并联须各自串栅极电阻、对称布线}。}
